% Save as: module8-problem-set.tex
% ============================================================================
% Problem Set 8: Environmental, Resource & Conservation
% Microeconomic Theory for Experimentalists & Applied Microeconomists
% Built from templates/problem_set_template.tex. Compile with pdflatex.
% ============================================================================
\documentclass[11pt]{article}
\usepackage{amsmath, amssymb, amsthm}
\usepackage[margin=1in]{geometry}
\usepackage{enumitem}

\newtheorem{question}{Question}

\title{Problem Set 8: Environmental, Resource \& Conservation\\
\large Microeconomic Theory for Experimentalists \& Applied Microeconomists}
\author{Due: \textbf{[date --- before Lecture 1 of Module 9]}}
\date{}

\begin{document}
\maketitle

\noindent\emph{Ground rules: collaboration on ideas is encouraged; write-ups
are individual. You may use AI tools for parts flagged with
$\langle$AI$\rangle$ and must attach the transcript. Show all calculations.}

\medskip

\noindent\textbf{Weights:} Part A --- 30 points. Part B --- 45 points.
Part C --- 25 points.

\bigskip

% ---------------------------------------------------------------------------
\section*{Part A: Theory --- Guided Calculations (30 points)}

\begin{question}[The commons, with arithmetic --- 15 points]
Four users share a resource pool that starts at 40 units. Each round,
each user requests an extraction of 0--10 units; if total requests
exceed the pool, it is split proportionally and exhausted. Whatever
remains after extraction grows by a factor of 1.5, capped at 40. The
game lasts six rounds (all players know this).

\begin{enumerate}[label=(\alph*)]
  \item \textbf{(4 pts)} Find the largest total per-round extraction
        $T$ that keeps the pool at 40 forever, and the corresponding
        symmetric per-user extraction (round to the nearest feasible
        integer). Compute each user's total earnings over six rounds on
        this sustainable path, and compare to the earnings from an
        all-grab in round one (everyone requests 10).
  \item \textbf{(4 pts)} The knife: consider a two-round version.
        Three users extract the sustainable 3 per round throughout. A
        fourth deviates, requesting 10 in round one and 3 in round two.
        Compute the deviator's and a cooperator's two-round earnings
        (track the pool level carefully), and state in one sentence why
        this makes the situation a social dilemma --- note where the
        deviation's cost appears, and why the full six-round horizon is
        what completes the dilemma.
  \item \textbf{(3 pts)} State the standard-theory prediction for this
        game played by rational money-maximizers without communication,
        and the prediction for the same game with a 30-second
        non-binding chat before each round. One sentence each.
  \item \textbf{(4 pts)} Ostrom, Walker \& Gardner (1992) found that
        face-to-face communication raises cooperation dramatically and
        durably. In one paragraph: why is this a puzzle for the theory
        in (c), and name the three candidate mechanisms from lecture
        (preferences, beliefs, identity) with one sentence on what each
        claims the talk is doing.
\end{enumerate}
\end{question}

\begin{question}[PES, inframarginality, and the auction --- 15 points]
A payment-for-ecosystem-services program pays forest owners \$20 per
hectare per year to keep forest standing. It enrolls 1{,}000 hectares.
Without payments, 10\% of these hectares would be deforested each year;
with payments, 5\% are. Each deforested hectare releases 100 tons of
CO$_2$.

\begin{enumerate}[label=(\alph*)]
  \item \textbf{(4 pts)} Compute the program's annual cost, the
        hectares of deforestation averted per year, the tons of CO$_2$
        averted, and the cost per ton averted.
  \item \textbf{(4 pts)} Consider a second program with the same terms
        but different enrollment: 500 of its 1{,}000 enrolled hectares
        belong to owners who would never have deforested (the payments
        to them are inframarginal); on its at-risk 500 hectares,
        deforestation falls from 10\% to 5\%. Compute this program's
        cost per ton averted, and state in one sentence why asking
        applicants ``did you plan to cut?'' cannot fix this.
  \item \textbf{(4 pts)} Explain how a procurement auction (Jack 2013)
        addresses inframarginality: what do bids reveal, who wins
        contracts under a budget, and why does self-selection do the
        targeting that screening questions cannot? Note one honest
        limitation of the auction (e.g.\ what happens if owners
        collude, or if opportunity cost and environmental value are
        negatively correlated).
  \item \textbf{(3 pts)} Much of the program's carbon benefit arrives
        near 2100. In three sentences a policymaker would accept: why does
        the choice between a 1\% and a 7\% discount rate change the
        program's present value by roughly two orders of magnitude,
        why is that choice normative rather than technical, and what
        should an honest cost-benefit report do about it?
\end{enumerate}
\end{question}

% ---------------------------------------------------------------------------
\section*{Part B: Experimental Design (45 points)}

\begin{question}[A conservation procurement auction, in the field]
Design a field experiment comparing a \emph{procurement auction} for
conservation contracts against a \emph{flat-rate} PES program of equal
budget, in a smallholder forest landscape. Specify:
\begin{enumerate}[label=(\alph*)]
  \item \textbf{(6 pts)} Setting and pool: the landscape, the
        landholders, the conservation contract's terms (what the owner
        must do, for how long), and why this setting makes the
        targeting question first-order.
  \item \textbf{(12 pts)} Treatments: (i) flat-rate offer arm, (ii)
        procurement-auction arm (specify the auction format --- sealed
        bid, discriminatory or uniform price --- and justify), and
        (iii) a pure control arm. Randomize at what level, and why?
        State precisely which comparisons identify (1) the program
        effect on deforestation, (2) the targeting gain from the
        auction (cost per averted hectare, auction vs.\ flat), and
        (3) what the bid distribution itself tells you.
  \item \textbf{(9 pts)} Measurement: compliance monitoring (satellite
        + ground checks), opportunity-cost proxies at baseline, and ---
        non-negotiable in this domain --- \textbf{leakage}: how you
        detect deforestation displaced to unenrolled plots or
        neighboring villages, and how the randomization level in (b)
        supports leakage measurement.
  \item \textbf{(9 pts)} Hypotheses and power: anchor the program
        effect on Jayachandran et al.\ (deforestation roughly halved)
        and state the cluster-randomized sample (villages per arm,
        landholders per village) needed, noting the design effect from
        clustering; state which comparison (program effect or
        auction-vs-flat targeting gain) drives the sample size and
        why.
  \item \textbf{(9 pts)} The biggest threat to YOUR design and your
        answer to it. (Candidates: leakage across arm boundaries;
        auction bids contaminated by anchoring on the flat rate if
        arms share information; monitoring-induced behavior change.)
\end{enumerate}
\end{question}

% ---------------------------------------------------------------------------
\section*{Part C: Silicon Pilot Interpretation $\langle$AI$\rangle$
(25 points)}

\begin{question}[The villagers who never broke a promise]
You run \texttt{module8-simulation-lab.ipynb}. With communication,
your silicon villagers' extraction drops to almost exactly the
sustainable level and stays there; the pre-round messages are
articulate mutual pledges; \textbf{every agent honors every pledge in
all rounds}, and there is no defection in the known final round. Human
groups with communication improve dramatically but imperfectly:
pledge-breaking occurs, cooperation erodes under temptation, and
known-endgame defection is common. Your lab partner reports:
``communication works even better than Ostrom found.''

\begin{enumerate}[label=(\alph*)]
  \item \textbf{(8 pts)} Explain why perfect pledge-keeping and the
        absence of endgame defection are \emph{divergences} from human
        behavior rather than a stronger version of Ostrom's result,
        and give two mechanisms (at least one LLM-specific) that could
        produce them.
  \item \textbf{(9 pts)} Design the diagnostic battery: your design
        must include (i) a \emph{temptation spike} (e.g.\ the final
        round's extraction value is doubled, and agents are told) and
        (ii) a \emph{message-ablation} arm (replace the agents' real
        messages with content-free filler while keeping the
        communication structure), and may add others (anonymity
        framing; a mixed group with one profit-focused persona). State
        the predicted result pattern under each mechanism from (a).
        Run it if API budget allows and report; otherwise state
        predictions precisely.
  \item \textbf{(8 pts)} A colleague designing a human CPR experiment
        with a communication treatment asks what your sim can offer.
        State one thing it CAN legitimately inform (and why), one
        thing it CANNOT (and why), and one concrete design change to
        the human experiment that your silicon results --- read
        correctly --- would justify.
\end{enumerate}
Attach your AI transcripts and, if you ran code, the results CSV.
\end{question}

\end{document}
